% === II. RELATED WORK ========================================================
\section{Related Work and Position}
\label{sec:related}

\subsection{On-track and data-driven identification}

The Magic Formula~\cite{pacejka1987sae,pacejka2012book} is the model identified
here. Classical characterisation obtains its coefficients from steady-state
manoeuvres in a controlled space \cite{ontrack2025,amz2020}: clean data at high
logistical cost. Direct fitting to measured contact-patch forces needs a
force-measuring hub or a simulator's own wheel telemetry; here it is used only
as a \emph{reference} against which the on-track identifier is scored, never
inside the identification loop.

On-track methods trade excitation quality for operational simplicity. Nonlinear
least squares on lap data is the classical baseline and is brittle under noise
\cite{ontrack2025}. The recursive alternative is a dual or joint filter carrying
states and parameters together \cite{wenzel2006dekf,sierra2006cornering}, and it
is subject to the same excitation condition this paper is about: a filter run on
a smooth racing line has as little information about the tire's peak as a batch
fit of the same data, and merely expresses that as a parameter covariance that
never shrinks rather than as a coefficient on a bound. Dedicated friction
estimators make the excitation dependence explicit
\cite{rajamani2012wheelfriction,ahn2013friction,acosta2017friction,khaleghian2017survey},
and Hsu \emph{et al.}~\cite{hsu2010trail} is the closest prior art to the warm
start of Section~\ref{subsec:friction}: they recover the friction limit below
the peak from the \emph{curvature} of the force response, read through the
pneumatic trail in steering torque. Our estimator uses the same observation on
the axle force curve reconstructed from the IMU and odometry, with no
steering-torque sensor; it is then gated, applied to both axles from the axle
the manoeuvre identifies, and pinned for the co-identification loop, which is a
property of what that loop can identify (Section~\ref{subsec:bcd}) rather than
of the estimator. Gaussian-process approaches learn the mismatch of a nominal
model \cite{hewing2018cautious,hewing2020cautious,kabzan2019lbmpc} at a
continuing computational cost; neural residuals replace the Gaussian process
with a cheaper regressor at the cost of generalisation guarantees.

Deep Dynamics~\cite{deepdynamics2024} is the closest work on the
physics-constraint axis: it estimates Pacejka coefficients inside a neural
network and adds a \emph{Physics Guard} layer clamping them into nominal ranges.
Our admissibility gates share the motivation but differ in placement. We report
as a substantive finding that per-coefficient box constraints are
insufficient: the degenerate fit that motivated this work satisfied every
box constraint while five of eight coefficients sat exactly on a bound, and a
later fit satisfied every box constraint while being pairwise oversteering and
open-loop unstable above \SI{14.5}{\meter\per\second}. We therefore gate on
derived quantities at the model-to-controller interface and treat a coefficient
resting on a bound as a diagnostic of non-identifiability rather than a
successful clamp. The hybrid residual-plus-virtual-data structure
of~\cite{ontrack2025} is the pipeline we extend; the idea it rests on, to query
the learned model only where it interpolates and extract physically
parameterised coefficients from it, is retained unchanged, and what we correct
is the mechanism deciding \emph{where} the learned model is queried.

\subsection{Overlap with concurrent work}

Wu \emph{et al.}~\cite{visionsysid2026} extend the same baseline concurrently
and independently, introducing a vision-based friction prior, an S4 structured
state-space residual in place of the MLP, and Nelder--Mead for the iterative
Pacejka extraction. The pipeline described here independently contains an S4D
temporal residual~\cite{gu2022s4d,gu2022s4}, a friction warm-start, and
Nelder--Mead and differential-evolution solver options, so we claim none of the
three as contributions. Three differences are technical rather than
chronological. Our friction prior is proprioceptive, so it needs no camera and
no textured-surface assumption. It is read from the \emph{curvature} of the
axle force curve by a gated brush-model
fit~\cite{fiala1954,acosta2017friction} rather than from utilised friction,
which on a gentle lap measures \num{0.044} median on this vehicle and is only a
lower bound (Section~\ref{subsec:frictionfloor}). And it is \emph{pinned}
rather than used as an initialisation, because Section~\ref{subsec:bcd} shows
the rollout fit cannot separate the peak factor from the cornering stiffness at
any excitation, so a prior on $D$ the fit is free to move is overwritten within
one cycle. Neither prior work is the fix: the degeneracy is located in the
rollout excitation, not the initialisation, and to our knowledge the
reachable-friction bound of Section~\ref{subsec:excitation} and its consequence
for the virtual-data step have not been stated before.

\subsection{Simulation platform and downstream consumer}

CARLA~\cite{carla2017} supplies the sensor and scenario infrastructure used
here. Its native dynamics come from the underlying PhysX wheeled-vehicle
model~\cite{physxvehicle}, a saturating friction-and-stiffness tire law
(Section~\ref{subsec:planttire}) rather than an empirical Magic Formula fit,
which is why this study treats the simulator as a well-instrumented \emph{plant}
to be identified and not as a source of ground-truth Pacejka coefficients.
R-CARLA~\cite{rcarla2025} makes the dynamics model interchangeable, and
Sagmeister \emph{et al.}~\cite{simfidelity2026} find that a controller's
sensitivity to model fidelity is largest near the acceleration limits, where
this study operates; CARLA has also hosted complete Formula Student Driverless
stacks with ROS~2 interfaces \cite{carlaros2025}. Scaled
platforms, F1TENTH~\cite{f1tenth2020} and the ForzaETH Race
Stack~\cite{forzaeth2024}, provide the software context of the baseline;
full-scale Formula Student Driverless systems~\cite{amz2020} the vehicle class
targeted here. The identified model is consumed by a nonlinear MPC path-tracking
controller built on \texttt{acados}~\cite{acados2022} with partial-condensing
HPIPM~\cite{hpipm2020} \cite{liniger2015,kabzan2019lbmpc}, tracking a
minimum-curvature raceline~\cite{heilmeier2020}. We claim no contribution to the
controller; it appears only as the consumer that defines what an admissible
identification is (Section~\ref{sec:integration}).
